% This document was generated by an LLM.
% Content fragment; compile intro-analysis-self-check.tex.
% Coverage: Hunter (2012), Chapter 1, printed pp. 1--10.
% Corrections to source slips:
% - Example 1.4 concerns nonempty finite sets.
% - Example 1.26 uses n >= 2, avoiding the reversed endpoints at n = 1.
% - Example 1.38 uses 1/(n+1), so every term belongs to (0,1).
% - Proposition 1.41 requires K to be nonempty.

\section{The Real Numbers}

\subsection{Completeness of \texorpdfstring{\(\R\)}{R}}
\textit{Source: \S1.1, pp.~1--3.}

\subsubsection*{Suprema and infima}
\begin{enumerate}
  % Definition 1.1 and the following equivalent boundedness condition.
  \item Define an upper bound and a lower bound of a set \(A\).
    What does it mean for \(A\) to be bounded above, bounded below, or bounded?
    State an equivalent definition of boundedness using absolute values.

  % Example 1.2.
  \item Give a set that is bounded below but unbounded above.
    Identify all its lower bounds and explain why it has no upper bound.

  % Definition 1.3.
  \item Define the supremum and infimum of a set. In each definition,
    distinguish the bound condition from the leastness or greatestness condition.

  % Definition 1.3; approximation property used in the proof of Proposition 1.41.
  \item Let \(A\subseteq\R\) be nonempty and let \(s\in\R\) be an
    upper bound of \(A\). State the \(\epsilon\)-approximation property
    equivalent to \(s=\sup A\), with all quantifiers.
    State the analogous characterization of \(m=\inf A\) for a lower
    bound \(m\in\R\).

  % Paragraph following Definition 1.3.
  \item Define the maximum and minimum of a set.
    How do they differ from the supremum and infimum?

  % Example 1.4; nonemptiness made explicit.
  \item What can be said about the boundedness, supremum, infimum,
    maximum, and minimum of a nonempty finite set of real numbers?

  % Example 1.5.
  \item Find the supremum and infimum of
    \(A=\{1/n:n\in\N\}\), and decide whether it has a maximum and a minimum.
    Verify the least-upper-bound and greatest-lower-bound conditions.

  % Example 1.6.
  \item Give an example of each of the following:
    a bounded nonempty set with neither a maximum nor a minimum;
    one with both; and one with a maximum but no minimum.
    Find the supremum and infimum of each.

  % Theorem 1.7.
  \item State the least-upper-bound formulation of completeness of \(\R\)
    (Theorem 1.7), including all hypotheses.
    Why can neither nonemptiness nor boundedness above be omitted?

  % Paragraph following Theorem 1.7.
  \item For \(-A=\{-x:x\in A\}\), explain the relationship between
    \(\inf A\) and \(\sup(-A)\), including when these real numbers exist.
    Deduce the greatest-lower-bound version of completeness.

  % Example 1.8.
  \item Give a nonempty subset of \(\Q\) that is bounded above by a rational
    number but has no supremum in \(\Q\). Explain why no rational upper bound
    is least. Does the same set have a supremum in \(\R\)?
\end{enumerate}

\subsubsection*{Cauchy sequences}
\begin{enumerate}[resume]
  % Definition 1.9.
  \item State the definition of a Cauchy sequence, with the quantifiers
    in order. How does it differ from the definition of convergence
    to a specified limit?

  % Theorem 1.10 and the surrounding discussion.
  \item State the Cauchy characterization of convergence for real sequences
    (Theorem 1.10). Which implication expresses completeness?

  % Example 1.11; an explicit example is requested in place of a given sequence.
  \item Give a Cauchy sequence of rational numbers that has no limit in
    \(\Q\). Explain why it is Cauchy and why it cannot converge to a
    rational number. What happens when it is regarded as a sequence in \(\R\)?
\end{enumerate}

\subsection{Open sets}
\textit{Source: \S1.2, pp.~3--5.}

\begin{enumerate}[resume]
  % Definition 1.12.
  \item State the definition of an open subset of \(\R\) using intervals
    of radius \(\delta\). May \(\delta\) depend on the point?

  % Definition 1.13.
  \item Define an interior point and a boundary point of a set.
    Must either kind of point belong to the set?

  % Paragraph following Definition 1.13.
  \item Characterize open sets in terms of interior points.

  % Example 1.14.
  \item Given \(x\in(a,b)\), with \(a<b\) finite, choose an explicit
    \(\delta>0\) such that \((x-\delta,x+\delta)\subseteq(a,b)\).
    How would you choose \(\delta\) for \(x\in(a,\infty)\) or
    \(x\in(-\infty,b)\)?

  % Definition 1.13 applied to the chapter's interval examples.
  \item Find all interior and boundary points of each of
    \((0,1)\), \([0,1]\), and \((0,1]\).

  % Unnumbered results following Example 1.14.
  \item State the closure properties of open sets under unions and
    intersections. Specify which collections may be arbitrary and which
    must be finite.

  % Example 1.15.
  \item Give a countable family of open intervals whose intersection is
    not open. Compute the intersection and explain why it is not open.

  % Unnumbered interval decomposition result following Example 1.14.
  \item State the interval decomposition result for open subsets of
    \(\R\). What can be required of the intervals and their number?

  % Introductory paragraph of Section 1.2.
  \item What does it mean for a property to be a topological property?
    Name two examples mentioned in the notes.

  % Definition 1.16.
  \item Define a neighborhood of \(x\) and a \(\delta\)-neighborhood of
    \(x\), using the convention in these notes. Must a neighborhood
    be open? Must it contain \(x\)?

  % Example 1.17.
  \item Give a neighborhood of a point that is not an open set.
    For \(a<b\), determine exactly which points of \([a,b]\) have
    \([a,b]\) as a neighborhood.

  % Definition 1.18.
  \item Restate the definition of an open set in terms of neighborhoods,
    and explain why it agrees with the interval definition.

  % Definition 1.19.
  \item Define what it means for \(B\subseteq A\) to be relatively open
    in \(A\). Which set in the definition must be open in \(\R\)?

  % Example 1.20.
  \item Give two subsets of \([0,1]\) that are open in \([0,1]\) but
    not open in \(\R\). For each, exhibit an open set \(U\subseteq\R\)
    whose intersection with \([0,1]\) is the given subset.

  % Definition 1.21.
  \item Define a relative neighborhood of \(x\in A\).
    Express the definition using a neighborhood in \(\R\).

  % Proposition 1.22.
  \item State the relative-neighborhood characterization of relative
    openness (Proposition 1.22).

  % Proof of Proposition 1.22.
  \item Prove Proposition 1.22 in both directions. In the direction that
    starts with a relative neighborhood at each point, construct an open
    set \(U\subseteq\R\) such that \(B=A\cap U\).
    Why is it useful to replace the given neighborhoods by open neighborhoods?
\end{enumerate}

\subsection{Closed sets}
\textit{Source: \S1.3, pp.~5--6.}

\begin{enumerate}[resume]
  % Definition 1.23.
  \item Define a closed subset of \(\R\) using its complement.
    In which ambient set is the complement taken?

  % Definition 1.24.
  \item Define a sequentially closed subset of \(\R\).
    Where must the sequence lie, and where is its convergence considered?

  % Unnumbered equivalence following Definition 1.24.
  \item State the relationship between closedness and sequential
    closedness for subsets of \(\R\).

  % Example 1.25.
  \item Verify that \([0,1]\) is closed in two ways: using its complement
    and using convergent sequences. Which property of limits is needed
    for the second argument?

  % Paragraph following Example 1.25.
  \item State the closure properties of closed sets under unions and
    intersections, specifying the permitted sizes of the collections.
    Explain how they follow from the corresponding facts about open sets.

  % Example 1.26; start at n = 2 to avoid reversed endpoints.
  \item Give a countable family of closed intervals whose union is not
    closed. Compute the union and explain why it is not closed.

  % Paragraph preceding Example 1.27.
  \item Which subsets of \(\R\) are both open and closed?
    Can a subset of \(\R\) be neither open nor closed?

  % Example 1.27.
  \item Give a set that is neither open nor closed. Show the failure of
    openness using neighborhoods and the failure of closedness using
    a convergent sequence.
\end{enumerate}

\subsection{Accumulation points and isolated points}
\textit{Source: \S1.4, pp.~6--7.}

\begin{enumerate}[resume]
  % Definition 1.28.
  \item Define an accumulation point of a set using intervals.
    What other names are used for an accumulation point?
    Why is the requirement that the nearby point differ from the given
    point essential?

  % Paragraph preceding Definition 1.29.
  \item Why must every neighborhood of an accumulation point of \(A\)
    contain infinitely many points of \(A\), rather than just one other point?

  % Definition 1.29.
  \item State the sequential characterization of an accumulation point.
    What restriction on the terms of the sequence must be included?
    Must those terms be pairwise distinct in this characterization?

  % Comparison of Definitions 1.13 and 1.28.
  \item Must every boundary point of a set be an accumulation point?
    Must every accumulation point be a boundary point?
    Give a simple counterexample to each implication that fails.

  % Explanation linking Definitions 1.28 and 1.29.
  \item Explain how to pass from the interval definition of an
    accumulation point to the sequential characterization and back.

  % Unnumbered closedness characterization following Definition 1.29.
  \item State the characterization of closed sets in terms of their
    accumulation points.

  % Example 1.30.
  \item Give an infinite subset of \(\R\) with no accumulation points
    in \(\R\). Explain why it has none.

  % Example 1.31.
  \item Find all accumulation points of \(A=\{1/n:n\in\N\}\).
    Which of them belong to \(A\)? Explain why there are no others.

  % Example 1.32.
  \item For \(a<b\), find all accumulation points of \((a,b)\)
    and of \([a,b]\). Which accumulation points belong to each set?

  % Example 1.33.
  \item For \(a<c<b\), find all accumulation points of
    \((a,c)\cup(c,b)\). What effect does removing \(c\) have on the
    set of accumulation points?

  % Definition 1.34 and the following dichotomy.
  \item Define an isolated point of a set. Must it belong to the set?
    State the dichotomy between isolated and accumulation points for
    points that belong to a set. Does the same dichotomy hold for every
    point of \(\R\)?

  % Example 1.35.
  \item Show that every point of \(\{1/n:n\in\N\}\) is isolated.
    For each \(n\), give an explicit positive radius that isolates \(1/n\),
    treating \(n=1\) separately if necessary.

  % Example 1.36.
  \item Can an interval have an isolated point? Distinguish intervals
    with more than one point from singleton intervals.
\end{enumerate}

\subsection{Compact sets}
\textit{Source: \S1.5, pp.~7--10.}

\subsubsection*{Sequential compactness}
\begin{enumerate}[resume]
  % Definition 1.37.
  \item Define sequential compactness. Which sequences are quantified
    over, and where must the subsequential limit belong?

  % Example 1.38; 1/(n+1) fixes the source's first term outside (0,1).
  \item Give a bounded set that is not sequentially compact.
    Exhibit a sequence in the set with no subsequence converging to a
    point of the set. Check that every term belongs to the set.

  % Example 1.39.
  \item Give a closed set that is not sequentially compact, and exhibit
    a sequence that demonstrates this failure.

  % Theorem 1.40.
  \item State the closed-and-bounded characterization of sequential
    compactness in \(\R\) (Theorem 1.40).

  % First paragraph of the proof of Theorem 1.40.
  \item Prove that a sequentially compact set is closed by considering
    a convergent sequence in the set. Why must the limit of a convergent
    subsequence agree with the limit of the original sequence?

  % Second paragraph of the proof of Theorem 1.40.
  \item Prove that a sequentially compact set is bounded.
    If the set were unbounded, how would you choose a sequence that
    rules out every convergent subsequence?

  % Last paragraph of the proof of Theorem 1.40.
  \item State the Bolzano--Weierstrass theorem for real sequences.
    Use it to prove that a closed, bounded subset of \(\R\) is
    sequentially compact. Where is each hypothesis used?

  % Proposition 1.41; nonemptiness is missing from the source statement.
  \item State the maximum-and-minimum result for sequentially compact
    subsets of \(\R\) (Proposition 1.41), including the necessary
    nonemptiness assumption. Why is that assumption needed?

  % Proof of Proposition 1.41.
  \item Prove that a nonempty sequentially compact set contains its
    supremum and infimum. Explain how the definition of supremum gives
    a suitable sequence, and where completeness and closedness enter.
\end{enumerate}

\subsubsection*{Open covers and compactness}
\begin{enumerate}[resume]
  % Unnumbered definitions preceding Definition 1.42.
  \item Define an open cover of a set and a finite subcover of a given
    open cover. Must the sets in a cover be contained in the set being
    covered? Must a subcover use sets from the original cover?

  % Definition 1.42.
  \item Define compactness using open covers. Is it enough to find
    one open cover that has a finite subcover?

  % Example 1.43.
  \item Construct an open cover of \(\N\) with no finite subcover.
    Explain why an arbitrary finite subcollection fails to cover \(\N\).

  % Final sentence of Example 1.43.
  \item Use the cover \(\{(-n,n):n\in\N\}\) to prove directly from
    the open-cover definition that every compact subset of \(\R\)
    is bounded.

  % Example 1.44.
  \item Construct an open cover of \((0,1)\) with no finite subcover.
    Verify that it covers every point and that any finite subcollection
    misses some points near an endpoint.

  % Examples 1.44--1.45; reproduce the notes' particular cover as an application.
  \item For integers \(n\geq0\), let
    \[
      I_n=\bigl(2^{-n}-2^{-(n+1)},\;2^{-n}+2^{-(n+1)}\bigr).
    \]
    Find \(\bigcup_{n=0}^{\infty}I_n\) and
    \(\bigcup_{n=0}^{N}I_n\). Why does this family cover \((0,1)\)
    but have no finite subcover of \((0,1)\)?

  % Example 1.45.
  \item For the same intervals \(I_n\), add \(J=(-\delta,\delta)\),
    where \(\delta>0\). Show that this gives an open cover of \([0,1]\)
    with a finite subcover by specifying a sufficient condition on \(N\)
    for \(J,I_0,\ldots,I_N\) to cover \([0,1]\).
    Why does verifying this particular cover alone not prove compactness?

  % Theorem 1.46 and its following discussion.
  \item State the closed-and-bounded characterization of compact subsets
    of \(\R\) (Theorem 1.46) and the interval version of the Heine--Borel
    theorem quoted in the notes.

  % Final paragraph of Chapter 1.
  \item What is the relationship between compactness and sequential
    compactness for subsets of \(\R\)? Explain how Theorems 1.40 and 1.46
    establish it.

  % Synthesis of the chapter's examples and characterizations.
  \item Decide which of the following sets are compact, and justify
    each answer:
    \[
      \varnothing,\quad \{2\},\quad [0,1],\quad (0,1],\quad
      \N,\quad \{1/n:n\in\N\},\quad \{0\}\cup\{1/n:n\in\N\}.
    \]
\end{enumerate}
